\chapter{Conclusions}
\label{chap:conclusion}

This thesis set out to build a systematic CFD screening methodology for swirl-recovery guide vanes downstream of the rim-driven thruster in the Store2Hydro booster-RPT concept, and to use it to map how the main design variables govern the trade-off between swirl reduction and hydraulic loss. Eighteen guide-vane configurations were simulated in the integrated RDT-GV-draft-tube domain, spanning two chord-control strategies, six shroud-chord levels and three vane counts, with a no-vane case as the reference.

The first research question asked how the guide-vane parameters influence swirl reduction, residual swirl distribution and total-pressure loss. The screening separated the matrix cleanly into two families. The constant-chord family gave a higher mean system head increase (\SI{3.70}{\meter} against \SI{3.49}{\meter}) and a lower mean total head loss, while the constant-solidity family reached the lowest residual flow angle once the chord was long enough. The guide-vane passage carried 58 to \SI{87}{\percent} of the head-loss budget. Useful swirl recovery required a shroud chord of at least \SI{400}{\milli\meter}, and no short-chord configuration reached the Pareto front.

The second and fourth questions concerned which metrics discriminate between designs and how robust the ranking is across the duct. The mass-averaged absolute flow angle and the swirl number at the cone end carried the comparison, supported by the radial angle profile and the Q-criterion strong-core fraction. For the Pareto candidates the bulk angle was representative of the radial structure, so a low bulk value did not hide inner-span weakness. Family-level and large-separation trends held across the several sampling planes and across the mesh correction, while close same-family rankings did not.

The third question asked whether the best design directly downstream of the vanes was also the best at the bend exit. The \SI{90}{\degree} bend amplified the swirl it received, so weak-recovery configurations were penalised through the bend while the strong-recovery candidates were not. The bend head loss collapsed by an order of magnitude, from \SI{0.37}{\meter} in the no-vane reference to below \SI{0.06}{\meter}, as soon as any vane was present. This collapse was binary and followed from the presence of a vane rather than from which vane, so the bend did not reorder the strong-recovery candidates.

A mesh-interface error at the RDT hub was found during post-processing and corrected on seven anchor cases. The correction confirmed the family-level trends and the Pareto status of \texttt{CS\_500\_8} and \texttt{CC\_500\_8}, while showing that close same-family rankings are case-sensitive. The recommendation is therefore \texttt{CS\_500\_8} where RPT inlet quality dominates the design priority, and \texttt{CC\_500\_8} where head delivery dominates, with \texttt{CC\_400\_8} as the highest-priority follow-up rerun. The transient analysis indicated hydraulic pulsation amplitudes well below the one-percent rotor-stator-interaction threshold, and equal rotor and stator counts were avoided in the shortlist as a conservative structural choice.

The study remains a comparative screening at one operating point and one inlet wake. Cavitation, the coupling to the RPT and the structural behaviour of the vanes were not modelled, and seven of the eighteen cases were rerun on the corrected mesh. These items set the scope of the recommendation and motivate the future work in Chapter~\ref{chap:future_work}.
%DEAR HÅVARD 10.06: (new content + humanizer) Replaced the placeholder "how to write a conclusion" boilerplate with a real conclusion built from Ch.1 (the four research questions) and Ch.4 (results, Pareto, mesh correction, synthesis). I removed the misplaced "Background" section heading, since a conclusion should not open with a Background section. Followed the template note: no \cite here, past tense throughout. VKL voice: SI units, "indicated" not "proved", British spelling, no em-dashes. Numbers are drawn from Ch.4 tables; please verify each against your final corrected matrix. (Note: the prior master-retting marker "CHANGE 2" that sat above this section flagged the same placeholder; it is now addressed.)

\clearpage % Keep this at the end of the chapter